\section{The separating family}\label{sec:construction}
Let $N\ge2$, $M=N^3$, and $K=2M$. Points are $p=(x,y)\in[N]\times[2N^2]$. Rows are affine lines $\ell=(a,b)\in[N]\times[2N^2]$. Put $F_{\ell p}=-1$ exactly when $y=ax+b$. A monotone flip matrix $A$ changes any subset of these $-1$ entries to $+1$ and leaves all other entries unchanged. In the random family the choices are independent fair bits. Write $H_A$ for its row class; \cref{fig:template} shows $N=2$.

\begin{figure}[t]
\centering
\begin{tikzpicture}[x=.235cm,y=.235cm,font=\tiny]
\begin{scope}[xshift=0.0cm]
\fill[white] (0,0) rectangle (16,16);
\fill[oiBlue!75] (1,15) rectangle (2,16);
\fill[oiBlue!75] (10,15) rectangle (11,16);
\fill[oiBlue!75] (2,14) rectangle (3,15);
\fill[oiBlue!75] (11,14) rectangle (12,15);
\fill[oiBlue!75] (3,13) rectangle (4,14);
\fill[oiBlue!75] (12,13) rectangle (13,14);
\fill[oiBlue!75] (4,12) rectangle (5,13);
\fill[oiBlue!75] (13,12) rectangle (14,13);
\fill[oiBlue!75] (5,11) rectangle (6,12);
\fill[oiBlue!75] (14,11) rectangle (15,12);
\fill[oiBlue!75] (6,10) rectangle (7,11);
\fill[oiBlue!75] (15,10) rectangle (16,11);
\fill[oiBlue!75] (7,9) rectangle (8,10);
\fill[oiBlue!75] (2,7) rectangle (3,8);
\fill[oiBlue!75] (12,7) rectangle (13,8);
\fill[oiBlue!75] (3,6) rectangle (4,7);
\fill[oiBlue!75] (13,6) rectangle (14,7);
\fill[oiBlue!75] (4,5) rectangle (5,6);
\fill[oiBlue!75] (14,5) rectangle (15,6);
\fill[oiBlue!75] (5,4) rectangle (6,5);
\fill[oiBlue!75] (15,4) rectangle (16,5);
\fill[oiBlue!75] (6,3) rectangle (7,4);
\fill[oiBlue!75] (7,2) rectangle (8,3);
\draw[step=1,black!20,line width=.15pt] (0,0) grid (16,16);\draw[black!70] (0,0) rectangle (16,16);
\node[anchor=east] at (-.3,15.5) {$(1,1)$};
\node[anchor=east] at (-.3,14.5) {$(1,2)$};
\node[anchor=east] at (-.3,13.5) {$(1,3)$};
\node[anchor=east] at (-.3,12.5) {$(1,4)$};
\node[anchor=east] at (-.3,11.5) {$(1,5)$};
\node[anchor=east] at (-.3,10.5) {$(1,6)$};
\node[anchor=east] at (-.3,9.5) {$(1,7)$};
\node[anchor=east] at (-.3,8.5) {$(1,8)$};
\node[anchor=east] at (-.3,7.5) {$(2,1)$};
\node[anchor=east] at (-.3,6.5) {$(2,2)$};
\node[anchor=east] at (-.3,5.5) {$(2,3)$};
\node[anchor=east] at (-.3,4.5) {$(2,4)$};
\node[anchor=east] at (-.3,3.5) {$(2,5)$};
\node[anchor=east] at (-.3,2.5) {$(2,6)$};
\node[anchor=east] at (-.3,1.5) {$(2,7)$};
\node[anchor=east] at (-.3,0.5) {$(2,8)$};
\node[rotate=60,anchor=east] at (0.5,-.25) {$(1,1)$};
\node[rotate=60,anchor=east] at (1.5,-.25) {$(1,2)$};
\node[rotate=60,anchor=east] at (2.5,-.25) {$(1,3)$};
\node[rotate=60,anchor=east] at (3.5,-.25) {$(1,4)$};
\node[rotate=60,anchor=east] at (4.5,-.25) {$(1,5)$};
\node[rotate=60,anchor=east] at (5.5,-.25) {$(1,6)$};
\node[rotate=60,anchor=east] at (6.5,-.25) {$(1,7)$};
\node[rotate=60,anchor=east] at (7.5,-.25) {$(1,8)$};
\node[rotate=60,anchor=east] at (8.5,-.25) {$(2,1)$};
\node[rotate=60,anchor=east] at (9.5,-.25) {$(2,2)$};
\node[rotate=60,anchor=east] at (10.5,-.25) {$(2,3)$};
\node[rotate=60,anchor=east] at (11.5,-.25) {$(2,4)$};
\node[rotate=60,anchor=east] at (12.5,-.25) {$(2,5)$};
\node[rotate=60,anchor=east] at (13.5,-.25) {$(2,6)$};
\node[rotate=60,anchor=east] at (14.5,-.25) {$(2,7)$};
\node[rotate=60,anchor=east] at (15.5,-.25) {$(2,8)$};
\node[font=\small,anchor=south] at (8,17) {Template $F$};
\node[rotate=90,font=\scriptsize] at (-4.5,8) {lines $(a,b)$};\node[font=\scriptsize] at (8,-3.9) {points $(x,y)$};\end{scope}
\begin{scope}[xshift=5.64cm]
\fill[white] (0,0) rectangle (16,16);
\fill[oiBlue!75] (1,15) rectangle (2,16);
\draw[oiOrange,very thick] (10.22,15.22)--(10.78,15.78);\draw[oiOrange,very thick] (10.22,15.78)--(10.78,15.22);
\fill[oiBlue!75] (2,14) rectangle (3,15);
\fill[oiBlue!75] (11,14) rectangle (12,15);
\draw[oiOrange,very thick] (3.22,13.22)--(3.7800000000000002,13.78);\draw[oiOrange,very thick] (3.22,13.78)--(3.7800000000000002,13.22);
\draw[oiOrange,very thick] (12.22,13.22)--(12.78,13.78);\draw[oiOrange,very thick] (12.22,13.78)--(12.78,13.22);
\draw[oiOrange,very thick] (4.22,12.22)--(4.78,12.78);\draw[oiOrange,very thick] (4.22,12.78)--(4.78,12.22);
\fill[oiBlue!75] (13,12) rectangle (14,13);
\fill[oiBlue!75] (5,11) rectangle (6,12);
\fill[oiBlue!75] (14,11) rectangle (15,12);
\fill[oiBlue!75] (6,10) rectangle (7,11);
\fill[oiBlue!75] (15,10) rectangle (16,11);
\draw[oiOrange,very thick] (7.22,9.22)--(7.78,9.78);\draw[oiOrange,very thick] (7.22,9.78)--(7.78,9.22);
\draw[oiOrange,very thick] (2.22,7.22)--(2.7800000000000002,7.78);\draw[oiOrange,very thick] (2.22,7.78)--(2.7800000000000002,7.22);
\fill[oiBlue!75] (12,7) rectangle (13,8);
\draw[oiOrange,very thick] (3.22,6.22)--(3.7800000000000002,6.78);\draw[oiOrange,very thick] (3.22,6.78)--(3.7800000000000002,6.22);
\draw[oiOrange,very thick] (13.22,6.22)--(13.78,6.78);\draw[oiOrange,very thick] (13.22,6.78)--(13.78,6.22);
\fill[oiBlue!75] (4,5) rectangle (5,6);
\draw[oiOrange,very thick] (14.22,5.22)--(14.78,5.78);\draw[oiOrange,very thick] (14.22,5.78)--(14.78,5.22);
\fill[oiBlue!75] (5,4) rectangle (6,5);
\fill[oiBlue!75] (15,4) rectangle (16,5);
\draw[oiOrange,very thick] (6.22,3.22)--(6.78,3.7800000000000002);\draw[oiOrange,very thick] (6.22,3.7800000000000002)--(6.78,3.22);
\draw[oiOrange,very thick] (7.22,2.22)--(7.78,2.7800000000000002);\draw[oiOrange,very thick] (7.22,2.7800000000000002)--(7.78,2.22);
\draw[step=1,black!20,line width=.15pt] (0,0) grid (16,16);\draw[black!70] (0,0) rectangle (16,16);
\node[anchor=east] at (-.3,15.5) {$(1,1)$};
\node[anchor=east] at (-.3,14.5) {$(1,2)$};
\node[anchor=east] at (-.3,13.5) {$(1,3)$};
\node[anchor=east] at (-.3,12.5) {$(1,4)$};
\node[anchor=east] at (-.3,11.5) {$(1,5)$};
\node[anchor=east] at (-.3,10.5) {$(1,6)$};
\node[anchor=east] at (-.3,9.5) {$(1,7)$};
\node[anchor=east] at (-.3,8.5) {$(1,8)$};
\node[anchor=east] at (-.3,7.5) {$(2,1)$};
\node[anchor=east] at (-.3,6.5) {$(2,2)$};
\node[anchor=east] at (-.3,5.5) {$(2,3)$};
\node[anchor=east] at (-.3,4.5) {$(2,4)$};
\node[anchor=east] at (-.3,3.5) {$(2,5)$};
\node[anchor=east] at (-.3,2.5) {$(2,6)$};
\node[anchor=east] at (-.3,1.5) {$(2,7)$};
\node[anchor=east] at (-.3,0.5) {$(2,8)$};
\node[rotate=60,anchor=east] at (0.5,-.25) {$(1,1)$};
\node[rotate=60,anchor=east] at (1.5,-.25) {$(1,2)$};
\node[rotate=60,anchor=east] at (2.5,-.25) {$(1,3)$};
\node[rotate=60,anchor=east] at (3.5,-.25) {$(1,4)$};
\node[rotate=60,anchor=east] at (4.5,-.25) {$(1,5)$};
\node[rotate=60,anchor=east] at (5.5,-.25) {$(1,6)$};
\node[rotate=60,anchor=east] at (6.5,-.25) {$(1,7)$};
\node[rotate=60,anchor=east] at (7.5,-.25) {$(1,8)$};
\node[rotate=60,anchor=east] at (8.5,-.25) {$(2,1)$};
\node[rotate=60,anchor=east] at (9.5,-.25) {$(2,2)$};
\node[rotate=60,anchor=east] at (10.5,-.25) {$(2,3)$};
\node[rotate=60,anchor=east] at (11.5,-.25) {$(2,4)$};
\node[rotate=60,anchor=east] at (12.5,-.25) {$(2,5)$};
\node[rotate=60,anchor=east] at (13.5,-.25) {$(2,6)$};
\node[rotate=60,anchor=east] at (14.5,-.25) {$(2,7)$};
\node[rotate=60,anchor=east] at (15.5,-.25) {$(2,8)$};
\node[font=\small,anchor=south] at (8,17) {Monotone flip $A$};
\node[rotate=90,font=\scriptsize] at (-4.5,8) {lines $(a,b)$};\node[font=\scriptsize] at (8,-3.9) {points $(x,y)$};\end{scope}
\end{tikzpicture}
\caption{The construction at $N=2$. Blue cells are negative incidences $y=ax+b$; white cells are positive. The right panel changes a fixed subset of the negative cells to positive, marked by orange crosses. Rows and columns retain their geometric indices. Larger scales supply the counting separation in \cref{thm:construction}.}\label{fig:template}
\end{figure}



\begin{theorem}[Construction and ordinary dimension]\label{thm:construction}
Every monotone flip $A$ satisfies $\operatorname{VC}(A)\le2$ and $\rect(A)\ge2^{-15}$. The template has $\sr(F)\le6$ and exactly
\begin{equation}\label{eq:I}
 I=2N^4-\left(\frac{N(N+1)}2\right)^2\ge N^4
\end{equation}
independent flip positions. Set $B_0=\log_2(4eK)$ and $s_0=\lfloor I/(4KB_0)\rfloor$. A random $A$ satisfies
\[
 \Pr[\dc(H_A)\le s_0]\le2^{-I/2}.
\]
Every member also satisfies $\dc(H_A)\le2N+1$. In particular most matrices have $\dc(H_A)=\Omega(N/\log M)$. For every $0<\epsilon<1/4$, all members have a learner with $m=O_\epsilon(\log M)$ and $\tau=\Omega_\epsilon(1)$, by \cref{thm:main} with $\rho=2^{15}$. Independently of the sharp rectangle input, the self-contained cap proof in Appendix~\ref{app:oldgeometry} supplies $\rho=2^{18}$ and the same asymptotic learning bound.
\end{theorem}

The proof uses the strict factorization
\[
 \langle(x^2,xy,y^2,x,y,1),(a^2,-2a,1,2ab,-2b,b^2-1/2)\rangle
 =(ax+b-y)^2-1/2.
\]
Positive template rectangles survive every flip. A negative template rectangle has a singleton side, so partitioning the other side by its new labels loses at most half its mass. Independent flips then supply $I$ free signs, to which Warren's counting theorem applies. Appendix~\ref{app:family-proof} gives the complete proof. Exactly $N(N+1)/2$ template rows have no incidences; all remain positive.


The restriction in \eqref{eq:structural} cannot be omitted. Every sign matrix is a monotone flip of the all-negative, rank-one matrix. More quantitatively, let $W$ be a $k\times k$ Walsh matrix, $k$ a power of two. Its orthogonality gives $\|W\|_\op=\sqrt k$. A monochromatic rectangle of side sizes $a,b$ satisfies
\[
 ab=|\1_S^TW\1_T|\le\sqrt{k}\sqrt{ab},
\]
so its uniform product mass is at most $1/k$. For $k>32$, this contradicts the asserted unrestricted bound $2^{-(2\cdot1+3)}$.


\subsection{Probabilistic dimension}

\begin{theorem}[Approximate counting and exact probabilistic dimension]\label{thm:prob}
Fix $0\le\eta<1/2$ and write
\[
 h_2(\eta)=-\eta\log_2\eta-(1-\eta)\log_2(1-\eta),\quad
 c_\eta=1-h_2(\eta)>0,
 \quad s_\eta=\left\lfloor\frac{c_\eta I}{4KB_0}\right\rfloor,
\]
with $0\log 0=0$. For the random incidence-flip matrix,
\begin{equation}\label{eq:approx-prob}
 \Pr[\dc^\eta(H_A)<s_\eta]\le2^{-c_\eta I/2}.
\end{equation}

Consequently most $A$ have $\dc_{1/2}(H_A)=\Omega(N/\log^2M)$ and $\dc^\eta(H_A)=\Omega_\eta(N/\log M)$.
\end{theorem}

\begin{proof}
Let $J_\ell$ be the set of template incidences in row $\ell$, with $k_\ell=|J_\ell|$. For each nonempty $J_\ell$, take $D_\ell$ uniform on it. If a law on $d$-dimensional embeddings witnesses $\dc^\eta(H_A)\le d$, then
\[
 \E_\phi\sum_{\ell:k_\ell>0}\frac{k_\ell}{I}
       \min_w L_{D_\ell,h_\ell}(\sign\langle w,\phi\rangle)\le\eta.
\]
For a fixed embedding, only finitely many Boolean predictions are possible, so each infimum is attained as a minimum of finitely many loss values. There exists one embedding with the displayed weighted sum at most $\eta$. Otherwise its positive excess over $\eta$ would have strictly positive expectation. Choose one minimizing weight vector for each row. Their prediction matrix makes at most $\eta I$ mistakes on the incidence positions.

With $\sign(0)=+1$, this prediction matrix has strict sign-rank at most $d+1$, not necessarily at most $d$. Add the constant feature one. In each row add a strictly positive bias smaller than the absolute value of every negative score in that row, if there are negative scores. This retains all signs and turns every zero into a positive score. Finiteness makes the bias choice possible. Rows with no incidence may use any weight vector.

For any fixed sign matrix $B$, the number of possible incidence flip patterns within $\eta I$ mistakes is at most
\[
 \sum_{j=0}^{\lfloor\eta I\rfloor}\binom Ij\le2^{h_2(\eta)I}.
\]
For $0<\eta<1/2$, this follows from
\[
 1\ge\sum_{j\le\eta I}\binom Ij\eta^j(1-\eta)^{I-j}
 \ge2^{-h_2(\eta)I}\sum_{j\le\eta I}\binom Ij;
\]
the second inequality uses $j\le\eta I$ and $\eta/(1-\eta)<1$. For $\eta=0$ the sum is one. Combining with \eqref{eq:warren}, the probability that any rank-at-most-$s_\eta$ sign matrix makes at most $\eta I$ incidence errors is at most
\[
 2^{2Ks_\eta B_0+h_2(\eta)I-I}\le2^{-c_\eta I/2}.
\]
An embedding of dimension below $s_\eta$ would produce just such a strict sign matrix after adding the constant coordinate. If $s_\eta=0$, the event in question is empty. This proves \eqref{eq:approx-prob}.

\end{proof}
\begin{lemma}[From exact probabilistic to strict dimension]\label{lem:exact-prob}
For every finite nonempty class $H$ with $K=|H|$,
\begin{equation}\label{eq:exact-prob}
 \dc(H)\le(\lceil\log_2K\rceil+1)\dc_{1/2}(H)+1.
\end{equation}
\end{lemma}
\begin{proof}
Choose a full-support distribution on $X$. Zero loss under this distribution means correct labels at every point. Draw $t=\lceil\log_2K\rceil+1$ independent embeddings from a law witnessing $\dc_{1/2}(H)=d$. Each fixed target fails in all $t$ embeddings with probability at most $2^{-t}$. A union bound gives probability at most $K2^{-t}<1$ that some target fails everywhere. Hence one collection works for all targets. Concatenate its $t$ embeddings; each row selects a successful block. Adding one constant coordinate as above makes the representation strict. This proves \eqref{eq:exact-prob}. 
\end{proof}



\begin{corollary}[Negative answers to the dimension implications]\label{cor:negative-answer}
For every fixed $0<\epsilon<1/4$, none of (L), (P) and (Pn) holds universally with $\dc(H)$ or $\dc_{1/2}(H)$ on the left. For every fixed $C>0$, choose
\[
 0<\epsilon<\min\{1/4,1/(2C)\},\qquad\eta=C\epsilon.
\]
The same three conclusions fail with $\dc^{C\epsilon}(H)$ on the left, under each nondegenerate tie convention of \cref{lem:ties}. The counterexamples have proper deterministic, table-polynomial learners with $m=O_\epsilon(\log M)$ and $\tau=\Omega_\epsilon(1)$.
\end{corollary}
\begin{proof}
Fix the required constant accuracies and intersect the high-probability events of \cref{thm:construction,thm:prob}. Their failure probability is at most $2^{-I/2}+2^{-(1-h_2(\eta))I/2}<1$ for all sufficiently large $N$. Thus one matrix satisfies the two lower bounds simultaneously. Apply \cref{lem:exact-prob} for exact probabilistic dimension, and \cref{thm:caps,thm:proper} for the proper deterministic learner. Since $M=N^3$ and $n=\log_2(2M)$, the dimension lower bounds grow exponentially in $n$, up to powers of $\log M$, whereas $m=O_\epsilon(n)$ and $1/\tau=O_\epsilon(1)$. No fixed polynomial in the proposed variables can dominate these lower bounds.
\end{proof}
\begin{table}[t]
\centering\small
\begin{tabularx}{\linewidth}{@{}p{.30\linewidth}X@{}}
\toprule
Resource & Incidence flips on $2M$ points, $M=N^3$\\\midrule
VC dimension; rectangle ratio & $\mathrm{VC}\le2$; $\rect\ge2^{-15}$.\\
Strict dimension, most flips & $\Omega(M^{1/3}/\log M)\le\dc\le2M^{1/3}+1$.\\
Exact probabilistic dimension & $\dc_{1/2}=\Omega(M^{1/3}/\log^2M)$.\\
Averaged dimension & $\dc^\eta=\Omega_\eta(M^{1/3}/\log M)$ for fixed $\eta<1/2$.\\
Table-known learners & $m=O_\epsilon(\log M)$, $\tau=\Omega_\epsilon(1)$; proper and deterministic; table-polynomial time.\\
Key-known median learner & $m=O_\epsilon(n^2)$, $\tau=\Omega_\epsilon(1)$; deterministic improper for every key; proper for regular stars.\\
Conditional high dimension & $\dc\ge n^c/(30\log_2n)$ for almost every PRF key, each fixed $c$.\\
Refuted conclusions & (L), (P), (Pn) for the three nondegenerate dimensions; every prescribed degree for the corresponding succinct families.\\
\bottomrule
\end{tabularx}
\caption{The separating family's parameters. The probabilistic lower bounds concern each prescribed fixed $\eta<1/2$ and hold with exponentially small failure probability. A polynomial-time evaluator with a supplied key is distinguished from deterministic selection of a good key.}\label{tab:family}
\end{table}

